\begin{summarybox}
	\textbf{\textcolor{CSummary}{一句话核心}}

	余弦定理变形公式直接由三边长度计算内角余弦，并可通过最大角的余弦符号判定三角形形状（正锐、零直、负钝）。

	\medskip
	\textbf{\textcolor{CSummary}{使用条件}}
	\begin{itemize}[leftmargin=*, itemsep=0.5em, topsep=0.4em]
		\item \textbf{条件 1（已知三边）：} 必须知道三角形三条边长或比例，且满足三角形不等式。
		\item \textbf{条件 2（最大角对应）：} 判断形状时需先找出最大边（从而确定最大角），再计算该角余弦。
	\end{itemize}

	\medskip
	\textbf{\textcolor{CSummary}{关键提醒}}
	\begin{itemize}[leftmargin=*, itemsep=0.5em, topsep=0.4em]
		\item \textbf{易错点（公式配对）：} 注意余弦公式分子中对边平方放在被减位置，分母是两边乘积的2倍，不可混淆。
		\item \textbf{检查项（符号含义）：} 余弦值为正不保证三角形锐角，只有最大角余弦为正才能说明锐角；零为直角；负为钝角。
	\end{itemize}
\end{summarybox}
